% GENERATED FILE. Edit canonical lessons, then run npm run generate:books.
% canonical-source-hash: fnv1a64:1f4c10535b522e04
% canonical-lessons: TE-C296-kathinamaina, TE-C296-pratyeka, TE-C296-anukulamaina, TE-C296-upayogakaramaina, TE-C296-sthanika

\chapter{Hard, Special, Convenient, Useful, Local}
\label{ch:te-hard-special-convenient-useful-local}
\hlchaptermodality{296}

\begin{chapteropening}
\textbf{By the end of this chapter:} \emph{I can say hard, special, convenient, useful and local.}

Complete the last lesson of chapter 296: I can say hard, special, convenient, useful and local.
\end{chapteropening}

\section[\texorpdfstring{\te{కఠినమైన} (kaṭhinamaina)}{కఠినమైన (kaṭhinamaina)}]{\te{కఠినమైన} (kaṭhinamaina) — hard, difficult; strict, harsh}

\label{lesson:TE-C296-kathinamaina}



\begin{quote}
Before the new one: say the Telugu for for some reason, then the Telugu for somehow.
\end{quote}

\subsection*{You'll want to know: \te{కఠినమైన}}
\textbf{\te{కఠినమైన}} — \emph{kaṭhinamaina} — \textquotedblleft{}hard, difficult; strict, harsh\textquotedblright{}.

\begin{grammarlens}[title={What you've built}]
One more word for this chapter.
\end{grammarlens}

\subsection*{Your turn}
\begin{itemize}
\raggedright
  \item \emph{Say it:} \emph{kaṭhinamaina}
  \item \emph{Say it:} \emph{kaṭhinamaina}, once more
  \item \emph{Say it:} say \emph{kaṭhinamaina}
  \item \emph{Recall:} say the Telugu for for some reason, then the Telugu for somehow, then say \emph{kaṭhinamaina} again
\end{itemize}

\begin{tcolorbox}[breakable,colback=teal!4,colframe=teal!35!black,title={Before you move on}]
What does \te{కఠినమైన} mean? (\textquotedblleft{}Hard, difficult; strict, harsh\textquotedblright{}.) Say it once more.
\end{tcolorbox}
\section[\texorpdfstring{\te{ప్రత్యేక} (pratyēka)}{ప్రత్యేక (pratyēka)}]{\te{ప్రత్యేక} (pratyēka) — special, separate}

\label{lesson:TE-C296-pratyeka}



\begin{quote}
Before the new one: say the Telugu for somehow, then the Telugu for hard.
\end{quote}

\subsection*{You'll want to know: \te{ప్రత్యేక}}
\textbf{\te{ప్రత్యేక}} — \emph{pratyēka} — \textquotedblleft{}special, separate\textquotedblright{}.

\begin{grammarlens}[title={What you've built}]
One more word for this chapter.
\end{grammarlens}

\subsection*{Your turn}
\begin{itemize}
\raggedright
  \item \emph{Say it:} \emph{pratyēka}
  \item \emph{Say it:} \emph{pratyēka}, once more
  \item \emph{Say it:} say \emph{pratyēka}
  \item \emph{Recall:} say the Telugu for somehow, then the Telugu for hard, then say \emph{pratyēka} again
\end{itemize}

\begin{tcolorbox}[breakable,colback=teal!4,colframe=teal!35!black,title={Before you move on}]
What does \te{ప్రత్యేక} mean? (\textquotedblleft{}Special, separate\textquotedblright{}.) Say it once more.
\end{tcolorbox}
\section[\texorpdfstring{\te{అనుకూలమైన} (anukūlamaina)}{అనుకూలమైన (anukūlamaina)}]{\te{అనుకూలమైన} (anukūlamaina) — convenient}

\label{lesson:TE-C296-anukulamaina}



\begin{quote}
Before the new one: say the Telugu for hard, then the Telugu for special.
\end{quote}

\subsection*{You'll want to know: \te{అనుకూలమైన}}
\textbf{\te{అనుకూలమైన}} — \emph{anukūlamaina} — \textquotedblleft{}convenient\textquotedblright{}.

\begin{grammarlens}[title={What you've built}]
One more word for this chapter.
\end{grammarlens}

\subsection*{Your turn}
\begin{itemize}
\raggedright
  \item \emph{Say it:} \emph{anukūlamaina}
  \item \emph{Say it:} \emph{anukūlamaina}, once more
  \item \emph{Say it:} say \emph{anukūlamaina}
  \item \emph{Recall:} say the Telugu for hard, then the Telugu for special, then say \emph{anukūlamaina} again
\end{itemize}

\begin{tcolorbox}[breakable,colback=teal!4,colframe=teal!35!black,title={Before you move on}]
What does \te{అనుకూలమైన} mean? (\textquotedblleft{}Convenient\textquotedblright{}.) Say it once more.
\end{tcolorbox}
\section[\texorpdfstring{\te{ఉపయోగకరమైన} (upayōgakaramaina)}{ఉపయోగకరమైన (upayōgakaramaina)}]{\te{ఉపయోగకరమైన} (upayōgakaramaina) — useful}

\label{lesson:TE-C296-upayogakaramaina}



\begin{quote}
Before the new one: say the Telugu for special, then the Telugu for convenient.
\end{quote}

\subsection*{You'll want to know: \te{ఉపయోగకరమైన}}
\textbf{\te{ఉపయోగకరమైన}} — \emph{upayōgakaramaina} — \textquotedblleft{}useful\textquotedblright{}.

\begin{grammarlens}[title={What you've built}]
One more word for this chapter.
\end{grammarlens}

\subsection*{Your turn}
\begin{itemize}
\raggedright
  \item \emph{Say it:} \emph{upayōgakaramaina}
  \item \emph{Say it:} \emph{upayōgakaramaina}, once more
  \item \emph{Say it:} say \emph{upayōgakaramaina}
  \item \emph{Recall:} say the Telugu for special, then the Telugu for convenient, then say \emph{upayōgakaramaina} again
\end{itemize}

\begin{tcolorbox}[breakable,colback=teal!4,colframe=teal!35!black,title={Before you move on}]
What does \te{ఉపయోగకరమైన} mean? (\textquotedblleft{}Useful\textquotedblright{}.) Say it once more.
\end{tcolorbox}
\section[\texorpdfstring{\te{స్థానిక} (sthānika)}{స్థానిక (sthānika)}]{\te{స్థానిక} (sthānika) — local}

\label{lesson:TE-C296-sthanika}



\begin{quote}
Before the new one: say the Telugu for convenient, then the Telugu for useful.
\end{quote}

\subsection*{You'll want to know: \te{స్థానిక}}
\textbf{\te{స్థానిక}} — \emph{sthānika} — \textquotedblleft{}local\textquotedblright{}.

\begin{grammarlens}[title={What you've built}]
One more word for this chapter.
\end{grammarlens}

\subsection*{Your turn}
\begin{itemize}
\raggedright
  \item \emph{Say it:} \emph{sthānika}
  \item \emph{Say it:} \emph{sthānika}, once more
  \item \emph{Say it:} say \emph{sthānika}
  \item \emph{Recall:} say the Telugu for convenient, then the Telugu for useful, then say \emph{sthānika} again
\end{itemize}

\begin{tcolorbox}[breakable,colback=teal!4,colframe=teal!35!black,title={Before you move on}]
What does \te{స్థానిక} mean? (\textquotedblleft{}Local\textquotedblright{}.) Say it once more.
\end{tcolorbox}
